\documentclass[11pt]{article}
\usepackage[margin=1in]{geometry}
\usepackage{hyperref}
\usepackage{enumitem}
\title{Block Coding with Sphero}
\author{Piter Garcia}
\date{\today}
\begin{document}
\maketitle

\section*{Grade and Class Match}
\textbf{Grade:} Grade 3\\
\textbf{Likely blocks:} Day 3 · 9:50-10:40 · Lamanaco; Day 4 · 9:50-10:40 · Lallucci; Day 5 · 9:50-10:40 · Regelsberger; Day 5 · 2:15-3:05 · Tandoi

\section*{Lesson Focus}
Block-based robot programming with movement, aiming, loops, and safety expectations.

\section*{Learning Targets}
\begin{itemize}[leftmargin=*]
\item Students connect to the correct robot and identify its controls.
\item Students create a short movement program using a loop.
\item Students follow equipment and partner expectations.
\end{itemize}

\section*{Materials}
\begin{itemize}[leftmargin=*]
\item Sphero robots
\item Student devices
\item Floor path or challenge space
\end{itemize}

\section*{Suggested Lesson Flow}
\begin{enumerate}[leftmargin=*]
\item Open with the exact device, tool, or routine students need in order to start successfully.
\item Model one example task before students begin independent or partner work.
\item Give students time to build, code, design, or document their work while circulating for support.
\item Close with a short share-out, reflection, or debugging discussion.
\end{enumerate}

\section*{Source Notes}
This lesson packet is enriched with transcript-derived lesson notes, the school schedule roster, and official starting-point resources. See the paired Markdown guide for clickable links.

\bibliographystyle{plain}
\bibliography{block-coding-with-sphero}
\end{document}
